\documentclass[12pt,a4paper]{article}

% --- Преамбула ---
\usepackage[utf8]{inputenc}
\usepackage[T1]{fontenc}
\usepackage[russian]{babel}
\usepackage{amsmath,amssymb}
\usepackage{hyperref}
\usepackage[margin=2.5cm]{geometry}
\usepackage{enumitem}
\usepackage{setspace}
\usepackage{booktabs}
\onehalfspacing

\title{\textbf{Измерение ортогональности электрона через статистику двухфотонного рождения пар}\\
\large Практический протокол для экспериментаторов}
\author{}
\date{7 июля 2026}

\begin{document}
\maketitle

\section{Что мы измеряем}
Ортогональность электрона $O_e$ — фундаментальный параметр Теории Мод, определяющий, как собственная мода электрона повёрнута относительно внешних мод. Теория предсказывает два возможных значения:

\begin{itemize}
    \item Модель A: $O_e \approx 0.999455$ (электрон почти полностью ортогонален, как фотон).
    \item Модель B: $O_e = 0.5$ (электрон ровно посередине между фотоном и протоном).
\end{itemize}

Обе модели согласуются с известными массами и константами связи, но дают разные предсказания для тонкой структуры собственной моды электрона. Мы предлагаем эксперимент, который различит их.

\section{Теоретическая основа}
В Теории Мод электрон обладает собственной модой с двумя нулями (слабый диполь). Расстояние между этими нулями — $K_0$ — определяет площадь локального участка моды $S_e$:

\begin{equation}
    S_e = m_e \cdot O_e
\end{equation}

где $m_e = 0.511$~МэВ — масса электрона.

Если средняя амплитуда поперечного спектра на участке между нулями порядка 1 (в безразмерных единицах), то:

\begin{equation}
    K_0 \approx S_e
\end{equation}

Для двух моделей:

\begin{table}[h]
    \centering
    \caption{Предсказания двух моделей}
    \begin{tabular}{lccc}
        \toprule
        Модель & $O_e$ & $S_e$ & $K_0$ \\
        \midrule
        A & 0.999455 & 0.5108 & 0.51 \\
        B & 0.5 & 0.2555 & 0.26 \\
        \bottomrule
    \end{tabular}
    \label{tab:models}
\end{table}

Разница в $K_0$ — примерно в 2 раза.

\section{Как $K_0$ влияет на рождение пар}
При двухфотонном рождении электрон-позитронной пары два фотона должны создать интерференционную картину с двумя нулями на расстоянии $K_0$. Это требует определённой разности частот:

\begin{equation}
    |\nu_1 - \nu_2| = \frac{c}{2 \cdot K_0 \cdot \delta}
\end{equation}

где $c$ — скорость света, $\delta$ — фундаментальная длина (порядка планковской, $\sim 10^{-35}$~м).

Разность частот для двух моделей:

\begin{table}[h]
    \centering
    \caption{Требуемая разность частот}
    \begin{tabular}{lcc}
        \toprule
        Модель & $K_0$ & Разность частот \\
        \midrule
        A & 0.51 & $\Delta\nu_A \approx c / (1.02 \cdot \delta)$ \\
        B & 0.26 & $\Delta\nu_B \approx c / (0.52 \cdot \delta) = 2 \cdot \Delta\nu_A$ \\
        \bottomrule
    \end{tabular}
    \label{tab:frequencies}
\end{table}

Модель B требует вдвое большей разности частот.

\section{Экспериментальный протокол}

\subsection{Установка}
\begin{itemize}
    \item Два лазера с перестраиваемой частотой (диапазон: рентгеновский или гамма, энергии порядка $0.5$--$1$~МэВ на фотон).
    \item Вакуумная камера с системой фокусировки двух пучков в одну точку.
    \item Детектор электрон-позитронных пар (калориметр + трековая система).
    \item Система точного контроля разности частот и относительной фазы двух лазеров.
\end{itemize}

\subsection{Процедура}
\begin{enumerate}[label=\arabic*.]
    \item Установить суммарную энергию двух фотонов $E_{\text{total}} = h\nu_1 + h\nu_2 = 1.022$~МэВ (порог рождения пары).
    \item Зафиксировать $\nu_1$ и начать сканировать $\nu_2$, меняя разность частот $\Delta\nu = |\nu_1 - \nu_2|$.
    \item Для каждого значения $\Delta\nu$ измерить выход пар (число зарегистрированных электрон-позитронных событий в единицу времени).
    \item Построить график: выход пар как функция $\Delta\nu$.
    \item Найти $\Delta\nu_{\text{opt}}$ — значение, при котором выход пар максимален.
\end{enumerate}

\subsection{Ожидаемый результат}
\begin{itemize}
    \item Если $\Delta\nu_{\text{opt}}$ соответствует $K_0 \approx 0.51$ → подтверждается модель A ($O_e \approx 0.999455$).
    \item Если $\Delta\nu_{\text{opt}}$ соответствует $K_0 \approx 0.26$ → подтверждается модель B ($O_e = 0.5$).
    \item Если пик не наблюдается или $K_0$ имеет другое значение → обе модели отвергаются, теория требует коррекции.
\end{itemize}

\section{Оценка требуемой точности}
Фундаментальная длина $\delta$ неизвестна, но мы можем обойтись без неё, измеряя \textbf{отношение} $K_0$ для электрона и позитрона. Поскольку позитрон имеет ту же массу и ту же ортогональность, его $K_0$ должно быть таким же.

Отношение $\Delta\nu_{\text{opt}} / E_{\text{total}}$ даст безразмерную величину, пропорциональную $1/K_0$. Измерив её с точностью $\sim 10\%$, мы сможем различить модели A и B, поскольку они отличаются в 2 раза.

\section{Альтернативный метод: однофотонный}
Если два лазера недоступны, можно использовать однофотонное рождение пар в поле ядра (процесс Бете-Гайтлера). В этом случае ядро играет роль второго фотона с нулевой частотой (статическое поле). Анализ распределения рождённых пар по углам и энергиям также чувствителен к $K_0$, хотя и слабее.

\section{Что это даст}
Измерение $O_e$ — это прямой тест Теории Мод. Если $O_e$ окажется равным $0.5$, это будет сильным аргументом в пользу логарифмической шкалы ортогональности и симметрии между фотоном, электроном и протоном. Если $O_e \approx 0.999455$ — подтвердится линейная модель с протоном как эталоном неортогональности. В любом случае, мы впервые измерим фундаментальный параметр, который стандартная физика не предсказывает и не измеряет.

\end{document}